\documentclass[11pt,a4paper]{article}
\usepackage[utf8]{inputenc}
\usepackage[T1]{fontenc}
\usepackage{lmodern}
\usepackage[french]{babel}
\usepackage{amsmath,amssymb,mathtools}
\usepackage{icomma}
\usepackage{xcolor}
\usepackage[most]{tcolorbox}
\usepackage{enumitem}
\usepackage[a4paper,margin=2.2cm,headheight=26pt,headsep=10pt]{geometry}
\usepackage{fancyhdr}
\usepackage[hidelinks]{hyperref}
\definecolor{primary}{RGB}{222,49,99}
\definecolor{primarydark}{RGB}{150,22,55}
\tcbset{boxbase/.style={enhanced,breakable,boxrule=0.9pt,arc=2.5pt,
  left=3mm,right=3mm,top=2.3mm,bottom=2.3mm,fonttitle=\bfseries,coltitle=white}}
\newtcolorbox[auto counter]{exo}[1][]{boxbase,colback=primary!4,colframe=primary,
  title={\textsc{Exercice}~\thetcbcounter\ifx\relax#1\relax\relax\else\textmd{ --- \itshape #1}\fi}}
\newcommand{\R}{\mathbb{R}}
\newcommand{\N}{\mathbb{N}}
\newcommand{\Z}{\mathbb{Z}}
\newcommand{\ds}{\displaystyle}
\pagestyle{fancy}\fancyhf{}
\fancyhead[L]{\small\textcolor{primarydark}{\textbf{Thème 1} — Puissances, racines et exposants}}
\fancyhead[R]{\small\textcolor{primarydark}{\textsc{Moyen}}}
\fancyfoot[C]{\small\thepage}
\renewcommand{\headrulewidth}{0.8pt}
\begin{document}

\begin{center}
{\LARGE\bfseries\color{primarydark} Niveau Moyen — Exercices}\\[2pt]
{\color{primary}\rule{0.45\textwidth}{1.2pt}}\\[3pt]
{\small\itshape On enchaîne maintenant plusieurs règles. Aucune calculatrice.}
\end{center}
\vspace{1mm}

\begin{exo}[combiner plusieurs règles]
Réduis chaque expression à une seule puissance $a^{k}$ (avec $a\in\R_0$) :
\begin{enumerate}[label=\alph*),leftmargin=2em,topsep=1pt,itemsep=1pt]
  \item $\ds\frac{a^{5}\times a^{-3}}{a^{2}}$ \qquad
  \item $\bigl(a^{2}\bigr)^{3}\times a^{-4}$ \qquad
  \item $\ds\frac{\bigl(a^{3}\bigr)^{2}}{a^{-1}}$ \qquad
  \item $\ds\Bigl(\frac{a^{-2}}{a^{3}}\Bigr)^{-1}$
\end{enumerate}
\end{exo}

\begin{exo}[racines et exposants fractionnaires]
Simplifie au maximum (donne le résultat sous forme $c^{k}$, $c>0$, $n\in\N_0$) :
\begin{enumerate}[label=\alph*),leftmargin=2em,topsep=1pt,itemsep=1pt]
  \item $\ds\frac{\sqrt{c^{3}}}{\sqrt[4]{c}}$ \qquad
  \item $\sqrt[3]{c}\times\sqrt[6]{c}$ \qquad
  \item $\ds\sqrt[n]{\sqrt[n]{c^{\,2n^{2}}}}$ \qquad
  \item $\ds\Bigl(\sqrt{c}\Bigr)^{6}\times c^{-2}$
\end{enumerate}
\end{exo}

\begin{exo}[bases opposées et quotients d'exposants]
Récris sous la forme d'un \emph{quotient} $\bigl(\tfrac{b}{a}\bigr)^{m}$ ou d'une puissance unique
($a,b\in\R_0$) :
\begin{enumerate}[label=\alph*),leftmargin=2em,topsep=1pt,itemsep=1pt]
  \item $a^{-m}\times b^{m}$ \qquad
  \item $\ds\frac{b^{3}}{a^{3}}$ \qquad
  \item $2^{-x}\times 3^{x}$ \qquad
  \item $\ds\frac{a^{4}b^{-4}}{1}$
\end{enumerate}
\end{exo}

\begin{exo}[ramener à une base commune avant de résoudre]
Résous dans $\R$ en écrivant d'abord les deux membres avec la \emph{même} base :
\begin{enumerate}[label=\alph*),leftmargin=2em,topsep=1pt,itemsep=1pt]
  \item $4^{x}=2^{6}$ \qquad
  \item $8^{x}=2^{x+4}$ \qquad
  \item $9^{x}=\ds\frac{1}{27}$ \qquad
  \item $\ds\Bigl(\frac{2^{2x+1}}{2^{x}}\Bigr)^{2}=8$
\end{enumerate}
\end{exo}

\begin{exo}[monotonie : comparer et résoudre une inéquation]
On rappelle : $a>1\Rightarrow x\mapsto a^{x}$ croissante ; $0<a<1\Rightarrow$ décroissante.
\begin{enumerate}[label=\alph*),leftmargin=2em,topsep=1pt,itemsep=1pt]
  \item Résous $2^{x}<2^{5}$.
  \item Résous $\ds\Bigl(\frac12\Bigr)^{x}>\Bigl(\frac12\Bigr)^{3}$ (attention au sens !).
  \item Range dans l'ordre croissant $2^{3}$, $2^{0}$, $2^{-1}$, $2^{1}$ sans les calculer.
  \item Résous $\ds\Bigl(\frac13\Bigr)^{2x}\leq \Bigl(\frac13\Bigr)^{6}$.
\end{enumerate}
\end{exo}

\end{document}
